\section{Threat Model and Scope}
\label{sec:threat}

% LOCKED (stage-1 decision record): semi-honest cloud + ML-capable attacker.
% This section MUST be precise — every "what breaks" claim is bounded by it.

\subsection{Adversary}
% Honest-but-curious cloud / GPU operator. Observes everything on the untrusted
% GPU: obfuscated activations, public model weights, access patterns, and stored
% ciphertexts. Has ML capability: can mount inversion, GRAM, and ISA AttnScore
% attacks. Follows the protocol (does not tamper with computation).
\needswork{formal-ish statement of the semi-honest adversary + observation set}

\subsection{Trust Anchors}
% Trusted: TEE seal, RATLS / TLS 1.3 channel, the client endpoint.
\needswork{enumerate trusted components and what each guarantees}

\subsection{Out of Scope}
% Explicitly OUT: active/malicious tampering (no integrity/verifiability claims);
% microarchitectural side-channels against the TEE (Spectre/Foreshadow-class).
% Stating these bounds pre-empts the "strawman adversary" objection.
\needswork{explicit out-of-scope list with one-line justification each}

\begin{table}[t]
  \centering
  \caption{Trust boundary assumed throughout the paper.}
  \label{tab:threat}
  \begin{tabular}{@{}lll@{}}
    \toprule
    Component & Status & Adversary capability \\
    \midrule
    Client endpoint        & Trusted   & --- \\
    TEE seal               & Trusted   & --- \\
    RATLS / TLS\,1.3       & Trusted   & --- \\
    Untrusted GPU          & Untrusted & sees obfuscated activations, weights \\
    Storage / vector DB    & Untrusted & sees ciphertexts, access patterns \\
    \midrule
    Active tampering       & \multicolumn{2}{l}{Out of scope} \\
    Microarch.\ side-chan. & \multicolumn{2}{l}{Out of scope} \\
    \bottomrule
  \end{tabular}
\end{table}
